\documentclass[11pt,a4paper]{article}
\usepackage[T1]{fontenc}
\usepackage{amsmath,amssymb,amsthm,mathtools,mathrsfs}
\usepackage[margin=23mm]{geometry}
\usepackage{longtable,booktabs,hyperref}
\hypersetup{colorlinks=true,linkcolor=blue,urlcolor=blue}
\allowdisplaybreaks
\setlength{\emergencystretch}{3em}
\newtheorem{theorem}{Theorem}[section]
\newtheorem{lemma}[theorem]{Lemma}
\newtheorem{proposition}[theorem]{Proposition}
\newtheorem{remark}[theorem]{Remark}
\newcommand{\C}{\mathbb C}
\newcommand{\R}{\mathbb R}
\newcommand{\Tr}{\operatorname{Tr}}
\newcommand{\diag}{\operatorname{diag}}
\newcommand{\im}{\operatorname{im}}
\newcommand{\rank}{\operatorname{rank}}
\newcommand{\Span}{\operatorname{span}}
\newcommand{\rem}{\operatorname{rem}}
\newcommand{\eps}{\varepsilon}
\title{The original Gamma return: joint full jets,\newline a closed boundary determinant, and mixed-word ranks}
\author{Split-Zero programme continuation}
\date{14 September 2026}
\begin{document}
\maketitle
\begin{abstract}
On the subsequence $k=4l+1$ we construct the exact Christoffel multiplication between the two original Gamma sources, calculate its joint arithmetic/boundary Schur complement, and return all four signs to $R_k^0-R_k^\sigma$. The unconditioned boundary determinant is evaluated in a closed product with a proof using two explicit adjoint polynomial maps. Its signed contribution has limit $(\log 2)/8$ after division by $q$ for a simple quartet; it is therefore not an error that can be discarded at the $q$ scale. The remaining mixed contribution is specified by matrices of size $l$, with a convergent finite signed trace enclosure and the original relation-valued section correction. We also give the actual coefficient formula for every mixed plus/minus word and a complete rank classification for the first nontrivial $k=5$, $(a,b)=(2,3)$ case. All arithmetic Grams remain those of the given source; no offcritical zero is asserted to exist.
\end{abstract}
\tableofcontents

\section{Scope, original input, and retained support}
The input is the actual counterfactual packet stipulated in the continuation request:
\begin{equation}
 h(s)=\prod_{\lambda\in\{\rho,\bar\rho,1-\rho,1-\bar\rho\}}(s-\lambda)^m,
 \quad\rho=\tfrac12+\delta+i\gamma,
 \quad0<\delta<\tfrac12,\quad\gamma>2.
\end{equation}
Its full multiplicities and arithmetic function are retained:
\begin{equation}
 g=2\xi,\quad v_h=g/h,\quad U=M_{j_hv_h},\quad
 w_h(y)=\frac{|v_h(\tfrac12+iy)|^2}{2\pi},\quad m_k=w_h^{*k}.
\end{equation}
Fix $k=4l+1$, $l\ge1$, put $c=k/2$, and set
\begin{equation}
 e_k=1+k(m-1),\quad q=(k+1)^2e_k,\quad
 \lambda_{ab}=c+(2a-k)\delta+i(2b-k)\gamma,
 \quad
 \chi(S)=\prod_{a,b=0}^k(S-\lambda_{ab})^{e_k}.
 \label{eq:packet}
\end{equation}
The ordered tensor fibre, all its multiplicities, and the weighted cyclic inclusion
\begin{equation}
 E=\C[S]/(\chi)\xrightarrow{\eta} E_h^{\otimes k},\qquad
 \eta[P]=\upsilon_h^{\otimes k}P\left(\sum_i s_i\right)1,
 \quad\upsilon_h=j_h(g/h),
 \label{eq:eta}
\end{equation}
are unchanged. In particular, a determinant of this rectangular inclusion is never inverted.

The scalar and support observations remain
\begin{equation}
 G(R)=\{\tau\}\sqcup R^\bullet,\quad e_R=0_R^\bullet,
 \quad(p_R,b_R)(\tau)=(0,0),\quad(p_R,b_R)(r^\bullet)=(r,1).
\end{equation}
The pair is injective. The structural map from the original pointed/preaddition base sends $\tau_0$ to $\tau$ and $1$ to $1_R^\bullet$. The arithmetic square retains its infinite integer target:
\begin{equation}
\begin{array}{ccc}
G(\mathbb Z)&\longrightarrow&G(\C)\\
p_{\mathbb Z}\downarrow&&\downarrow p_{\C}\\
\mathbb Z&\longrightarrow&\C.
\end{array}
\end{equation}
Every linear coefficient map constructed below is applied on each existing face and original outer leg. On a fixed present label $\omega$ its lift is $(\omega,x)\mapsto(\omega,Lx)$; external absence maps to external absence. On face inclusions $j_A^B$, zero insertion gives $L^B j_A^B=j_A^B L^A$. A represented zero, including an empty coefficient face at a present outer label, is not discarded. The existing intertype carrier maps and their nonunital sections retain their original types. These calculations concern their coefficient-linear arithmetic retracts; they do not assert that a linear section is a unital algebra homomorphism.

On the original distinguished tensor class the supplied same-class map is
\begin{equation}
 I_k(v_k)=\frac{(k(m-1))!}{((m-1)!)^k}v_h(\rho)^k\,u_\rho^{\otimes k}\ne0.
\end{equation}
Multiplication by the comparison polynomial below carries this eigenclass to $f(k\rho)I_k(v_k)$, with $f(k\rho)\ne0$; its explicitly constructed inverse restores it. On the surrounding primary module the complete triangular Leibniz matrix, not just this eigenvalue, is retained.

The original source remains
\begin{align}
 V&=\{\phi\in\mathcal S(\R)_{\rm ev}:\phi(0)=0,\ \int\phi=0\},\\
 \mathscr B&=\{F\in C^\infty(\R_{>0}):\sup_{x>0}x^b|D^jF(x)|<\infty\ (b\in\mathbb Z,j\ge0)\},\\
 D&=-x\partial_x,\quad \Theta\phi(x)=2\sum_{n\ge1}\phi(nx),\quad
 \mathcal MF(s)=\int_0^\infty F(x)x^s\,dx/x.
\end{align}
We use the existing Euler inverses only on their proved domains, $\mathcal MF_h=v_h$, and
\begin{equation}
 \mathcal V_{h,k}P=P(D_1+\cdots+D_k)F_h^{\otimes k},\qquad
 J^{(k)}\mathcal V_{h,k}=\eta\rem_\chi.
\end{equation}
All actual source norms still integrate $|P(c+iy)|^2m_k(y)dy$. The Gamma norms used below are comparison norms on this same polynomial presentation.

\paragraph{Inherited determinant input.}
DMR/QDC, as stated in the new request, give $|\delta_k^\sigma|\le2\Xi_{k,2q}$ and the asserted rate for $\Xi$. This note does not reattribute that theorem or replace its full determinant proof by diagonal estimates. The supplied proof-file inventory is recorded separately. When a requested proof body is absent from that inventory, its theorem is an explicitly inherited input, not a claim of independent verification. All results from Section~\ref{sec:gamma} through Section~\ref{sec:rank} below have their proofs here and do not assume the missing growing-family upper bound.

\section{The exact two-reference multiplication}\label{sec:gamma}
Retain
\begin{equation}
 \sigma(y)=\frac{|\Gamma(1/4+iy/2)|^2}{2\pi},\quad c_\sigma=\sqrt{2\pi},\quad
 c_a=2^{1-2a}\Gamma(2a),\quad
 r_k^0(y)=\frac{c_\sigma^k}{c_{k/4}}\frac{|\Gamma(k/4+iy/2)|^2}{2\pi}.
\end{equation}
Their masses are respectively $c_\sigma$ and $c_\sigma^k$.
Put $b_j=2j+\tfrac12$ and
\begin{equation}
 F_l^\pm(y)=\prod_{j=0}^{l-1}(y\pm ib_j),\quad
 f_l(S)=\prod_{j=0}^{l-1}(S-\alpha_j),\quad
 \alpha_j=c+b_j=2l+2j+1.
\end{equation}
Gamma recurrence gives the holomorphic identity
\begin{equation}
 \frac{\Gamma(l+1/4+iy/2)}{\Gamma(1/4+iy/2)}
 =(i/2)^l F_l^-(y).
\end{equation}
For real $y$, $F_l^- =\overline{F_l^+}$ and $F_l^+(y)=(-i)^lf_l(c+iy)$. Therefore
\begin{equation}
 \boxed{r_k^0(y)=C_k|f_l(c+iy)|^2\sigma(y),\qquad
 C_k=\frac{c_\sigma^k}{c_{k/4}}4^{-l}.}
 \label{eq:gamma}
\end{equation}
The plus polynomial has not been substituted in the holomorphic recurrence: conjugation occurs at the displayed real-line observation. The alternative monic polynomial with roots $c-b_j$ is related by $f_l=(-1)^l(f_l^-)^{\dagger_c}$, where $P^{\dagger_c}(S)=\overline{P(2c-\bar S)}$. Multiplication by $f_l/f_l^-$ is the corresponding modulus-one map between their source images on this line. Its arithmetic action is $M_{f_l}M_{f_l^-}^{-1}$, not the identity.

For $k=5$, $l=1$, the complete formula is
\begin{equation}
 f_1(S)=S-3,\qquad C_5=16\pi^2/3.
\end{equation}
On the domain \eqref{eq:packet}, every $\lambda_{ab}$ has nonzero imaginary part because $k$ is odd. Thus $f_l$ and $\chi$ are coprime. The factor $M_f:E\to E$ is an $E$-module automorphism, with inverse the unique Bezout remainder. It satisfies
\begin{equation}
 \eta M_f=M_{f(\sum s_i)}\eta.
\end{equation}
All full-unit factors in \eqref{eq:eta} commute with this multiplication and stay in that equation.

\subsection{Raw confluent jets and original remainders}
Let $J_N:\mathcal P_N\to E$ be the monic remainder map. In raw derivative coordinates let
\begin{equation}
 \mathcal J_\chi P=(P^{(r)}(\lambda_{ab}))_{a,b,\,0\le r<e_k}.
\end{equation}
The coordinate map $\mathcal J_\chi=\mathsf E_\chi J_N$ has entries
\begin{equation}
 (\mathsf E_\chi)_{(\lambda,r),n}
 =\begin{cases}\dfrac{n!}{(n-r)!}\lambda^{n-r}&n\ge r,\\0&n<r.\end{cases}
\end{equation}
Its determinant is
\begin{equation}
 \det\mathsf E_\chi=
 \left(\prod_{\lambda}\prod_{r=0}^{e_k-1}r!\right)
 \prod_{\lambda<\mu}(\mu-\lambda)^{e_k^2},
\end{equation}
for the stated lexicographic block order. Its inverse is the full Hermite interpolation map, not value-only interpolation. Multiplication by $f$ has jet entries
\begin{equation}
 (\mathsf F_\chi)_{r,s}=\binom rs f^{(r-s)}(\lambda)\quad(r\ge s)
\end{equation}
inside each block and zero between blocks. Hence
$M_f=\mathsf E_\chi^{-1}\mathsf F_\chi\mathsf E_\chi$ and
$\det M_f=\prod_\lambda f(\lambda)^{e_k}$.

\section{Joint full jets and the mixed Schur complement}\label{sec:schur}
Fix $N\ge q-1$ and put $M=N+l$. All sources embed in the common $\mathcal P_{2q+l}$.
Let $H_M^\sigma$ be its unchanged sigma Gram. Let
\begin{equation}
 Z_M:P\longmapsto(P(\alpha_0),\ldots,P(\alpha_{l-1})),\qquad
 \mathscr J_M=\binom{J_M}{Z_M}.
\end{equation}
Coprimality and monic division prove that $\mathscr J_M$ is onto, with kernel
$\chi f\mathcal P_{M-q-l}$. In particular its full joint covariance is positive definite:
\begin{equation}
 \mathscr J_M(H_M^\sigma)^{-1}\mathscr J_M^*
 =\begin{pmatrix}K_M&C_M\\C_M^*&A_M\end{pmatrix}.
 \label{eq:joint}
\end{equation}
Here $K_M=J_M(H_M^\sigma)^{-1}J_M^*$, $G_M^\sigma=K_M^{-1}$, and $A_M$ is the covariance of the additional real nodes. Define both positive Schur complements
\begin{equation}
 S_M=K_M-C_MA_M^{-1}C_M^*,\qquad
 L_M=A_M-C_M^*G_M^\sigma C_M.
 \label{eq:schurs}
\end{equation}
These formulas retain the actual cross block $C_M$.

\begin{theorem}[Original tensor-Gamma quotient]
The canonical tensor-Gamma quotient is
\begin{equation}
 \boxed{G_N^0=C_k M_f^*S_M^{-1}M_f.}\label{eq:quotient}
\end{equation}
Its inverse and its finite-rank return are
\begin{align}
 (G_N^0)^{-1}&=C_k^{-1}M_f^{-1}S_MM_f^{-*},\\
 G_N^0&=C_kM_f^*G_M^\sigma M_f
 +C_k(C_M^*G_M^\sigma M_f)^*L_M^{-1}(C_M^*G_M^\sigma M_f).
 \label{eq:woodbury}
\end{align}
\end{theorem}
\begin{proof}
Multiplication $\mathfrak m_f:\mathcal P_N\to\mathcal P_M$ is injective and has image $\ker Z_M$. Its source norm is exactly $C_k^{-1}$ times the original tensor-Gamma norm by \eqref{eq:gamma}. The arithmetic constraint $J_NP=v$ is transported to $J_M(fP)=M_fv$. Thus its minimum is the minimum sigma norm with joint constraint $(M_fv,0)$. Inverting \eqref{eq:joint} gives \eqref{eq:quotient}. Multiplication of the two block factors proves the Schur inverse identity in \eqref{eq:woodbury}; it uses $L_M$, not a covariance with $C_M$ erased.
\end{proof}
The representative itself is
\begin{equation}
 R_N^0=\operatorname{div}_f
 (H_M^\sigma)^{-1}\mathscr J_M^*
 \begin{pmatrix}K_M&C_M\\C_M^*&A_M\end{pmatrix}^{-1}
 \binom{M_f}{0}.
\end{equation}
The division is defined on $\ker Z_M=f\mathcal P_N$. Both $R_N^0$ and $R_N^\sigma$ satisfy $J_NR=1$. With $B_N$ the original multiplication-by-$\chi$ columns,
\begin{equation}
 R_N^0=R_N^\sigma-B_N(B_N^*H_N^0B_N)^{-1}B_N^*H_N^0R_N^\sigma.
 \label{eq:section}
\end{equation}
Thus their difference is the original relation, with its actual coefficient. The source Pythagoras identity retains
\begin{equation}
 (R_N^\sigma)^*H_N^0R_N^\sigma-G_N^0
 =(R_N^\sigma-R_N^0)^*H_N^0(R_N^\sigma-R_N^0).
\end{equation}
Applying $\mathcal V_{h,k}$ and the original fixed-order division by the factors $h(s_i)$ gives its theta primitive. The sign $(-1)^{i-1}$ in that primitive meets the same tensor-differential sign. At $N=q-1$ the relation space is zero; every expression containing its projector is the zero map.

Equation \eqref{eq:woodbury} is an application of the already supplied selected-mode/Schur update: its base term is explicitly retained, and the update has rank at most $l$. It is not a new abstract Woodbury principle.

\section{The exact signed return consists of small matrices}\label{sec:signed}
Let the four endpoints and signs be
\begin{equation}
 (N_0,N_1,N_2,N_3)=(q-1,q,2q-1,2q),\qquad
 (\eps_0,\eps_1,\eps_2,\eps_3)=(1,1,-1,-1).
\end{equation}
In the fixed monic sigma polynomial basis put
\begin{equation}
 U_N=[Jp_{N+1},\ldots,Jp_{N+l}],\quad
 \Omega_N=\diag(\omega_{N+1}^\sigma,\ldots,\omega_{N+l}^\sigma),\quad
 X_N=\Omega_N^{-1}U_N^*G_N^\sigma U_N.
\end{equation}
$X_N$ is positive self-adjoint in $\Omega_N$; its entries retain all confluent derivatives through $J$. Define the second, differently metrized endomorphism
\begin{equation}
 Y_N=A_{N+l}^{-1}C_{N+l}^*G_{N+l}^\sigma C_{N+l}.
\end{equation}
It is self-adjoint in $A_{N+l}$ and has spectrum in $[0,1)$, since $L_{N+l}>0$.
The original kernel update gives
\begin{equation}
 d_N:=\log\frac{V_N^\sigma}{V_{N+l}^\sigma}=\log\det(I+X_N),\qquad
 \beta_N:=\log\frac{\det A_{N+l}}{\det L_{N+l}}=-\log\det(I-Y_N).
\end{equation}
Taking the determinant of \eqref{eq:quotient} yields
\begin{equation}
 \log V_N^0=q\log C_k+2\log|\det M_f|+\log V_{N+l}^\sigma+\beta_N.
\end{equation}
The first two constants have not been removed before constructing these four metrics. Their signed sum is zero. In the exact notation of the incoming TWA return,
\begin{equation}
 \boxed{R_k^0-R_k^\sigma
 =\mathcal B_k^\sigma-\mathcal B_k^0
 =\sum_{r=0}^3\eps_r\{\log\det(I+X_{N_r})+\log\det(I-Y_{N_r})\}.}
 \label{eq:return}
\end{equation}
The matrices in \eqref{eq:return} have size $l=(k-1)/4$, while the retained arithmetic module still has dimension $q$.

\subsection{A positive determinant expression for the conditional block}
Let $A_n^{\chi\sigma}$ be the evaluation covariance at the same real nodes for polynomials of degree at most $n$ in the original measure $|\chi(c+iy)|^2\sigma(y)dy$. Orthogonal projection onto $\ker J_M=\chi\mathcal P_{M-q}$ proves
\begin{equation}
 \boxed{L_M=\diag(\chi(\alpha_j))\,A_{M-q}^{\chi\sigma}\,
 \diag(\overline{\chi(\alpha_j)}).}\label{eq:conditional}
\end{equation}
All $\chi(\alpha_j)$ are nonzero; more explicitly they are the positive products
\begin{equation}
 \chi(\alpha_j)=
 \prod_{a=0}^k\prod_{b=0}^{(k-1)/2}
 \left((\alpha_j-c-(2a-k)\delta)^2+(2b-k)^2\gamma^2\right)^{e_k}.
\end{equation}
The common multiplier in \eqref{eq:conditional} cancels only in the signed sum. At the first endpoint $N=q-1$ one can also evaluate the joint determinant directly:
\begin{equation}
 d_{q-1}-\beta_{q-1}
 =\log\frac{\left|\prod_j\chi(\alpha_j)\right|^2\Delta_f^2}
 {\left(\prod_{n=q}^{q+l-1}\omega_n^\sigma\right)\det A_{q+l-1}},
 \quad
 \Delta_f=2^{l(l-1)/2}\prod_{r=1}^{l-1}r!.
 \label{eq:first}
\end{equation}
Indeed the joint evaluation on the ordered basis $(1,\ldots,S^{q-1},\chi,\ldots,\chi S^{l-1})$ has determinant $(\prod_j\chi(\alpha_j))\Delta_f$, and the basis change has determinant one. No derivative factorial is lost: raw-jet coordinates introduce the same $\det\mathsf E_\chi$ on the two sides and their equality was given explicitly above.

\section{Closed evaluation of the free boundary determinant}\label{sec:closed}
The sigma monic polynomials satisfy, in $x=S-c$,
\begin{equation}
 p_0=1,\quad p_1=x,\quad
 p_{n+1}=xp_n+n(n-\tfrac12)p_{n-1},\qquad
 \omega_n^\sigma=c_\sigma n!(\tfrac12)_n.
 \label{eq:rec}
\end{equation}
All masses remain in $\omega_n^\sigma$. Write $\vartheta=1/2$. At the real node $\alpha_j$,
\begin{equation}
 p_n(\alpha_j)=(\vartheta)_n Q_j(n),\qquad
 Q_j(n)=\sum_{r=0}^j\binom jr\frac{2^r n^{\underline r}}{(\vartheta)_r}.
 \label{eq:eval}
\end{equation}
\begin{proof}
The exponential generating function derived from \eqref{eq:rec} solves
$(1-z^2)F'=(x+z/2)F$, $F(0)=1$. At $x=2j+1/2$ it is
$(1+z)^j(1-z)^{-j-1/2}$. Expanding the right side of \eqref{eq:eval}, using
$\sum_n(\vartheta)_n n^{\underline r}z^n/n!=(\vartheta)_r z^r(1-z)^{-\vartheta-r}$, gives the same series. Its leading coefficient as a polynomial in $n$ is $2^j/(\vartheta)_j$.
\end{proof}
It follows that
\begin{equation}
 (A_M)_{ij}=\frac1{c_\sigma}\sum_{n=0}^M
 \frac{(\vartheta)_n}{n!}Q_i(n)Q_j(n).
 \label{eq:discrete}
\end{equation}
This is an exact finite polynomial Gram, not an asymptotic kernel.

\subsection{Proof of its monic norms with both measures retained}
For any $\vartheta>0$ define the discrete weights $w_n=(\vartheta)_n/n!$, $0\le n\le M$, and the continuous measure
\begin{equation}
 K_Mt^{\vartheta-1}dt,\qquad K_M=(\vartheta)_{M+1}/M!.
\end{equation}
Their common mass is $(\vartheta)_{M+1}/(\vartheta M!)$. The joint measure
\begin{equation}
 K_M\binom Mn t^{n+\vartheta-1}(1-t)^{M-n}\,dt
\end{equation}
has exactly these marginals by the beta integral. It defines the adjoint polynomial maps
\begin{equation}
 (BP)(t)=\sum_{n=0}^M\binom Mn t^n(1-t)^{M-n}P(n),\qquad
 B^*(t^r)(n)=\frac{(n+\vartheta)_r}{(M+\vartheta+1)_r}.
\end{equation}
The adjunction is in the two displayed original measures, with their common mass retained.

Let $P_j$ be the monic continuous orthogonal polynomial. The formula
\begin{equation}
 P_j(t)=\frac{(-1)^j}{(j+\vartheta)_j}
 t^{1-\vartheta}\frac{d^j}{dt^j}
 \{t^{j+\vartheta-1}(1-t)^j\}
\end{equation}
has leading coefficient one. Integrating by parts $j$ times proves orthogonality and
\begin{equation}
 \|P_j\|^2=K_M\frac{(j!)^2}{(2j+\vartheta)(j+\vartheta)_j^2}.
\end{equation}
Every boundary term vanishes: at zero the last relevant derivative is $O(t^{\vartheta})$, and at one the factor $(1-t)^j$ supplies the required zeros.
Let $Q_j^{\rm mon}$ be the discrete monic polynomial. Adjointness and comparison of leading terms give
\begin{equation}
 BQ_j^{\rm mon}=M^{\underline j}P_j,\qquad
 B^*P_j=\frac{Q_j^{\rm mon}}{(M+\vartheta+1)_j}.
\end{equation}
Taking the same cross pairing in both measures proves its exact norm
\begin{equation}
 \nu_j(M)=
 \frac{(\vartheta)_{M+j+1}(j!)^2}
 {(M-j)!(2j+\vartheta)(j+\vartheta)_j^2}.
 \label{eq:discretenorm}
\end{equation}
This argument holds for every $0\le j\le M$ and supplies all determinant factors.

\begin{theorem}[Entire free boundary determinant]
For $M\ge l-1$ and $\vartheta=1/2$,
\begin{equation}
 \boxed{
 \det A_M=c_\sigma^{-l}\prod_{j=0}^{l-1}
 \frac{4^j(j!)^2(\vartheta)_{M+j+1}}
 {(\vartheta)_j^2(M-j)!(2j+\vartheta)(j+\vartheta)_j^2}.}
 \label{eq:detfree}
\end{equation}
\end{theorem}
\begin{proof}
In \eqref{eq:discrete}, changing the fixed monomial basis to $Q_j$ has determinant $\prod_j2^j/(\vartheta)_j$. The monic orthogonal-polynomial change has determinant one. Thus the Gram determinant is $c_\sigma^{-l}$ times the square of the first determinant times $\prod_j\nu_j(M)$. Substitute \eqref{eq:discretenorm}.
\end{proof}
In particular,
\begin{equation}
 \frac{\det A_{M+1}}{\det A_M}
 =\prod_{j=0}^{l-1}\frac{M+j+\vartheta+1}{M+1-j}.
 \label{eq:freeincrement}
\end{equation}

\subsection{Its signed growing-family contribution is evaluated}
Define
\begin{equation}
 C_{q,l}^{\rm free}
 =\log\frac{\det A_{2q+l-1}\det A_{2q+l}}
 {\det A_{q+l-1}\det A_{q+l}},\qquad B_l=l(l-\tfrac12).
\end{equation}
The expression is strictly positive. Formula \eqref{eq:freeincrement} evaluates it by finite products. It also gives the uniform finite bounds
\begin{equation}
 \boxed{
 2B_l\log\frac{2q+2l+1/2}{q+2l+1/2}
 \ \le C_{q,l}^{\rm free}\le\ 2B_l\log2.}
 \label{eq:freebound}
\end{equation}
\begin{proof}
For $t_j=2j+\vartheta$, each logarithmic increment is
$\log(1+t_j/(M+1-j))$. The inequalities $x/(1+x)\le\log(1+x)\le x$ apply with their positive denominators. Sum the two length-$q$ windows in \eqref{eq:freeincrement}. Writing their indices as $n=q,\ldots,2q-1$, all lower-bound denominators are at most $n+2l+\vartheta$, and all upper-bound denominators are at least $n+1$. Since $\sum_jt_j=B_l$, integration of the decreasing functions $1/(n+c)$ gives \eqref{eq:freebound}.
\end{proof}
The same bound gives the finite remainder
\begin{equation}
 0\le 2B_l\log2-C_{q,l}^{\rm free}
 \le \frac{B_l(2l+1/2)}q.
\end{equation}
Indeed the difference of the two logarithms is
$\log(1+(2l+1/2)/(2q+2l+1/2))$, at most $(2l+1/2)/(2q)$.
For simple quartets, $q=(4l+2)^2$, so the squeeze gives
\begin{equation}
 \boxed{\lim_{l\to\infty}\frac{C_{q,l}^{\rm free}}q=\frac{\log2}{8}.}
 \label{eq:qlimit}
\end{equation}
For each fixed $m>1$, $q=(4l+2)^2(1+(4l+1)(m-1))$, and the same ratio tends to zero. This is an evaluated component of the signed return, not a statement about the total mixed determinant.

Combining \eqref{eq:return} and \eqref{eq:conditional} gives the exact decomposition
\begin{equation}
 \boxed{
 R_k^0-R_k^\sigma=C_{q,l}^{\rm free}+
 \sum_{r=0}^3\eps_r
 \{d_{N_r}+\log\det A_{N_r+l-q}^{\chi\sigma}\}.}
 \label{eq:freeplusmixed}
\end{equation}
The source/relation-dependent summand in \eqref{eq:freeplusmixed} is retained. In particular \eqref{eq:qlimit} is not a limit for $R_k^0-R_k^\sigma$ itself.

\section{A finite signed enclosure and the exact return scale}\label{sec:bound}
Here is a finite enclosure for the entire expression \eqref{eq:return}, not just its free component. Use the common scalar
\begin{equation}
 s=2+\sum_{r=0}^3\Tr X_{N_r},\qquad
 Z_N=I-(I+X_N)/s.
\end{equation}
Each $Z_N$ is positive self-adjoint in $\Omega_N$ with spectrum in $[0,1)$; $Y_N$ already has that property in $A_{N+l}$. Retaining the common determinant factor first gives
\begin{equation}
 R_k^0-R_k^\sigma
 =-\sum_r\eps_r\sum_{n\ge1}
 \frac{\Tr Z_{N_r}^n+\Tr Y_{N_r}^n}{n}.
 \label{eq:series}
\end{equation}
The four $l\log s$ terms cancel before bounds are applied. There is no commutation assumption between different blocks.

For any one of these eight endomorphisms $T$, set
\begin{equation}
 b_p(T)=\sum_{n=1}^p\frac{\Tr T^n}{n},\quad a_p(T)=\Tr T^p,\quad
 u_p(T)=b_p(T)+\frac{b_{2p}(T)-b_p(T)}{1-a_p(T)}.
\end{equation}
For any $p$ with $a_p(T)<1$, the inherited positive-tail estimate gives
$b_p(T)\le-\log\det(I-T)\le u_p(T)$. To check its applicability here, diagonalize only in the stated positive metric: all eigenvalues are in $[0,1)$ by the proven Schur positivity. A finite such $p$ therefore exists for each given matrix. Termwise $r_{2p}(x)\le x^pr_p(x)$ proves the estimate without an additional arithmetic hypothesis.

Let $L_r=b_{p_r}(Z_{N_r})+b_{q_r}(Y_{N_r})$ and
$U_r=u_{p_r}(Z_{N_r})+u_{q_r}(Y_{N_r})$, using any admitted orders. Then the entire signed return satisfies
\begin{equation}
 \boxed{L_2+L_3-U_0-U_1\le R_k^0-R_k^\sigma
 \le U_2+U_3-L_0-L_1.}\label{eq:cert}
\end{equation}
Increasing the finite orders converges to the exact value. This is the existing PR29/STC trace certificate applied to the new small, fully specified joint-Schur matrices; it is not attributed as a new general logarithm inequality. All input errors must be included when numerical entries are used.

\subsection{Return to the original arithmetic criterion}
The inherited TWA equality and DMR/QDC bound give
\begin{equation}
 \mathcal B_k^{\rm ar}=\mathcal B_k^0+(R_k^0-R_k^\sigma)+\delta_k^\sigma,
 \qquad |\delta_k^\sigma|\le2\Xi_{k,2q}.
\end{equation}
If $[L_k,U_k]$ is the explicitly constructed interval \eqref{eq:cert}, the resulting arithmetic enclosure is the actual interval
\begin{equation}
 \boxed{\mathcal B_k^{\rm ar}\in
 [\mathcal B_k^0+L_k-2\Xi_{k,2q},\,
  \mathcal B_k^0+U_k+2\Xi_{k,2q}].}\label{eq:arithmeticinterval}
\end{equation}
The supplied original amplification identities also give, on either length-$q$ window,
\begin{equation}
 L_{h,k}^{2q}\le
 \frac{\omega_j^{\rm ar}}{\omega_i^{\rm ar}}\frac{V_i^{\rm ar}}{V_j^{\rm ar}},
 \quad L_{h,k}\ge\frac\delta2 kq,
\end{equation}
with the separate first-admissible-degree formula at $i=q-1$. Thus the original HC5 bound $4q\log(D_hk)$ is intersected with \eqref{eq:arithmeticinterval}; its original $D_h$ is not redefined here. The computed upper slack is
\begin{equation}
 \mathcal B_k^0+U_k+2\Xi_{k,2q}-4q\log(D_hk).
\end{equation}
This note has not proved that slack negative on all actual offcritical packets.

The error scale is explicit. A log-volume error bounded by $E_k$ changes a window exponential by factors between $e^{-E_k/(2q)}$ and $e^{E_k/(2q)}$. Thus $o(q)$ gives a factor tending to one; $o(q\log k)$ preserves the leading logarithmic threshold. The supplied DMR rate, with its original $p=8m-9/2$, gives only the available estimate
\begin{equation}
 \frac{\Xi_{k,2q}}q
 =O_h\left(k^{1/p}(\log(k+2))^{1-1/p}\right).
\end{equation}
This upper estimate does not imply either of the preceding smaller error scales. It also does not prove the actual error grows at that rate. The explicitly evaluated \eqref{eq:qlimit} shows why a $q$-scale reference term cannot be discarded even when it is $o(kq)$.

\section{Every mixed word and its actual observation}\label{sec:rank}
This section works on the reduced pair observations of the original full quartet, not on a new assertion that every original root is simple. All higher jets remain in the kernel of the reduced maps. Put
\begin{equation}
 c_\pm=\tfrac12\pm\delta,\qquad
 h_\pm(s)=(s-c_\pm)^2+\gamma^2.
\end{equation}
The reduced remainder from the full packet is onto each indicated value algebra. The reduced maps are not used as faithful observations of the full packet: for the original top local eigenvector $v_k$ at $k\rho$, every reduced value map kills it when $e_k>1$. When $e_k=1$, the word map to real-margin $a$ kills it unless $a=k$. These claims follow by reducing its local representative $(S-k\rho)^{e_k-1}$ and by the distinct real margins in \eqref{eq:packet}. The actual nonzero class remains in the specified kernel before that quotient. The source continues to use $v_h$, not the different literal divisor function $v_{h_\pm}=h/h_\pm\,v_h$.

Let $\Pi_\pm(u,0)$ be the convergent oriented pair period matrices in their monomial frames, and let $W$ be the retained invertible character-coordinate matrix of the cubic cover. Define
\begin{equation}
 E_\pm=\begin{pmatrix}1&c_\pm-i\gamma\\1&c_\pm+i\gamma\end{pmatrix},\quad
 U_\pm=\diag(v_h(c_\pm-i\gamma),v_h(c_\pm+i\gamma)),\quad
 T_\pm=W^{-1}\Pi_\pm E_\pm^{-1}U_\pm.
 \label{eq:pairmaps}
\end{equation}
The determinant of $E_\pm$ is $2i\gamma$. Both $U_\pm$ are invertible at the marked actual quartet, and every entry is retained. These maps go from original reduced jet values, through weighted polynomial coefficients, to actual character period coordinates.

For a word with $a$ plus and $b=k-a$ minus factors, use the ordered unscaled symmetric orbit vectors. Their coordinate Gram entries are
$\binom ai\binom bj$; the actual period and theta forms are still obtained by their displayed pullbacks, not assigned this coordinate Gram. Grouping an ordered word uses its explicit permutation and the top-cochain Koszul sign.

Write the two rows of $T_+$ as linear forms $A(X,Y),B(X,Y)$, and those of $T_-$ as $C(X,Y),D(X,Y)$. The selected rows are
\begin{equation}
 \mathcal I_{a,b}=\{(i,j):0\le i\le a,\ 0\le j\le b,\ k+i+j\equiv0\pmod3\}.
\end{equation}
For each such row expand the actual degree-$k$ form
\begin{equation}
 A^{a-i}B^iC^{b-j}D^j=\sum_{r=0}^k\mathsf C_{(i,j),r}X^{k-r}Y^r.
 \label{eq:rowpoly}
\end{equation}
These are explicit finite binomial sums in the entries of \eqref{eq:pairmaps}; no rank is assigned from the row count.
Put
\begin{equation}
 \lambda_r=ac_++bc_-+(2r-k)i\gamma,
 \qquad \mathsf E_{a,b}=(\lambda_r^n)_{0\le r,n\le k}.
\end{equation}
The roots are distinct, and
$\det\mathsf E_{a,b}=(2i\gamma)^{k(k+1)/2}\prod_{j=1}^k j!$.
The exact selected-row map on the cyclic reduced coefficient algebra is
\begin{equation}
 \boxed{\Lambda_{a,b}=\mathsf C\mathsf E_{a,b}.}\label{eq:wordmap}
\end{equation}
\begin{proof}
In an unscaled orbit sum, the coefficient of a fixed output word is the product of its row linear forms. Summing all input words with exactly $r$ upper-root entries extracts the coefficient of $X^{k-r}Y^r$. The sum operator on that input has eigenvalue $\lambda_r$. Evaluation of the cyclic polynomial gives precisely \eqref{eq:wordmap}. This also proves that no extra inverse binomial coefficient belongs in the formula.
\end{proof}
For arbitrary $a,b$ its exact rank is determined by this finite coefficient matrix, with no maximal-rank presumption. One exact algebraic implementation forms the coefficients of $\det(I+t\mathsf C^*H\mathsf C)$ with the retained positive row Gram $H$; the largest nonzero coefficient is the rank. A lexicographically first nonzero minor supplies a basis of the image and, by solving the retained pivot columns against the nonpivot columns, the complete kernel. Multiplication by $\mathsf E_{a,b}^{-1}$ returns that kernel to the original cyclic coefficient frame. These are finite coefficient calculations for the actual period entries, not an assumption that an unevaluated minor is nonzero.

\subsection{Completed first odd-subsequence mixed case: $k=5,a=2,b=3$}
The selected rows are $(0,1),(1,0),(1,3),(2,2)$. Their polynomials are
\begin{equation}
 A^2C^2D,\qquad ABC^3,\qquad ABD^3,\qquad B^2CD^2.
 \label{eq:fourrows}
\end{equation}
The invertible coordinate map $(X,Y)\mapsto(A,B)$ and its induced $\operatorname{Sym}^5$ map have their full inverse from $T_+^{-1}$. In those coordinates write
\begin{equation}
 C=\alpha A+\beta B,\quad D=\gamma_0 A+\delta_0 B,
 \qquad\begin{pmatrix}\alpha&\beta\\\gamma_0&\delta_0\end{pmatrix}
 =T_-T_+^{-1}.
\end{equation}
\begin{theorem}[True rank, including the exceptional locus]
The rank of \eqref{eq:fourrows} is two when $\beta=\gamma_0=0$, and four otherwise. The reduced cyclic kernel consequently has dimension four or two, respectively. Its full-quartet preimage has dimension $q-2$ or $q-4$, retaining all omitted roots and all higher jets.
\end{theorem}
\begin{proof}
When $C=\alpha A,D=\delta_0B$, the first two rows are nonzero multiples of $A^4B$ and the last two are multiples of $AB^4$, proving rank two.
If all four coefficients are nonzero, evaluate a linear dependence at $A=0$ to kill the fourth coefficient and at $B=0$ to kill the first. The remaining dependence is $AB(uC^3+vD^3)=0$, and independent $C,D$ force $u=v=0$.
It remains to check the possible zero coefficients of an invertible two-by-two matrix. If $\beta=0,\gamma_0\ne0$, evaluation at $B=0$ first kills the first row coefficient. After division by the common nonzero polynomial $AB$, evaluation at $D=0$ kills the second, and the coefficients of $A,B$ kill the other two. The case $\gamma_0=0,\beta\ne0$ is the reversed argument: evaluate at $A=0$, divide by $AB$, then evaluate at $C=0$. If $\alpha=0,\delta_0\ne0$, evaluation at $A=0$ kills the fourth coefficient; after division by $AB$, evaluation at $B=0$ kills the third and the remaining two are independent because $\gamma_0\ne0$. If $\delta_0=0,\alpha\ne0$, reverse that argument. Finally $\alpha=\delta_0=0$ gives four distinct monomials $A^3B^2,AB^4,A^4B,A^2B^3$. These cases exhaust every invertible matrix.
\end{proof}
For rank two the kernel in the transformed binary coefficient frame is explicitly the annihilator of the $A^4B$ and $AB^4$ coefficients. For rank four take the four rows in their stated order; any first nonzero four-column minor gives the exact two free-column kernel vectors by its matrix inverse. This returns through the explicit $\operatorname{Sym}^5T_+$ and $\mathsf E_{2,3}$ maps; no auxiliary metric changes the kernel.

\subsection{The two potentials' discrepancy is calculated}
The translation $s=c_\pm+z$ has Jacobian one and preserves the oriented contour differences. Thus
\begin{equation}
 \Pi_\pm E_\pm^{-1}=e^{C_\pm/u}\Pi_{\rm cen}E_0^{-1},\qquad
 C_\pm=c_\pm^3/3+c_\pm\gamma^2.
\end{equation}
Put $B_0=W^{-1}\Pi_{\rm cen}E_0^{-1}$ and $z_0=v_h(\rho)\ne0$. Reflection and conjugation of the original even-degree quartet give
\begin{equation}
 U_+=\diag(\bar z_0,z_0),\quad U_-=\diag(z_0,\bar z_0),\quad
 T_-T_+^{-1}=e^{(C_--C_+)/u}
 B_0\diag(z_0/\bar z_0,\bar z_0/z_0)B_0^{-1}.
\end{equation}
This is the exact discrepancy entering the rank theorem. For $B_0=\left(\begin{smallmatrix}a_0&b_0\\c_0&d_0\end{smallmatrix}\right)$ the offdiagonal entries are, respectively,
\begin{equation}
 \frac{a_0b_0(\bar z_0/z_0-z_0/\bar z_0)}{\det B_0},\qquad
 \frac{c_0d_0(z_0/\bar z_0-\bar z_0/z_0)}{\det B_0},
\end{equation}
with the same retained exponential scalar. Consequently the exceptional rank-two locus is exactly the vanishing of both these expressions. No phase condition on the actual $z_0$ is assumed.

\section{Joint boundary covariance and the same arithmetic class}
Let $p_{a,b}:E\twoheadrightarrow E_{a,b}$ be the actual reduced remainder. Its kernel retains every other primary block and the higher jets in the selected blocks. Its metric is the induced original quotient
\begin{equation}
 G_{a,b}=(p_{a,b}G^{-1}p_{a,b}^*)^{-1},\qquad
 r_{a,b}=G^{-1}p_{a,b}^*G_{a,b}.
\end{equation}
Then $p_{a,b}r_{a,b}=1$ and $1-r_{a,b}p_{a,b}$ projects onto its full kernel in the original $G$ metric.

Combine the existing observation $\Lambda$ and the calculated word map by
\begin{equation}
 \mathcal O=\binom{\Lambda}{\Lambda_{a,b}p_{a,b}}:
 E\twoheadrightarrow\im\mathcal O.
\end{equation}
On its actual image the covariance and minimum section are
\begin{equation}
 \Sigma=\mathcal O G^{-1}\mathcal O^*,\qquad
 s=G^{-1}\mathcal O^*\Sigma^{-1},\qquad
 \ker\mathcal O=\ker\Lambda\cap p_{a,b}^{-1}(\ker\Lambda_{a,b}).
\end{equation}
In the larger product target the covariance has all four blocks
\begin{equation}
 \begin{pmatrix}
 \Lambda G^{-1}\Lambda^*&\Lambda G^{-1}p_{a,b}^*\Lambda_{a,b}^*\\
 \Lambda_{a,b}p_{a,b}G^{-1}\Lambda^*&
 \Lambda_{a,b}p_{a,b}G^{-1}p_{a,b}^*\Lambda_{a,b}^*
 \end{pmatrix}.
\end{equation}
Its inverse is taken only on $\im\mathcal O$. The conditional Schur complement may have a kernel and that kernel is retained; no independent-observation assumption removes either cross term. Substituting \eqref{eq:quotient} or its inverse computes this covariance for the exact tensor-Gamma return. Substituting the original arithmetic $G$ computes it for the unchanged theta source.

At general roots $\zeta$ of $h-t$, if $v_h(\zeta+z)=z^\nu a(z)$ with $a$ a unit, multiplication on $\C[z]/z^r$ has image $(z^{\min(\nu,r)})$ and kernel $(z^{r-\min(\nu,r)})$. This is the explicit extension of the marked unit map used above; it is not an unproved invertibility claim away from the marked fibre. The source Fourier signs, oriented period contours, and nilpotent exponentials remain those of the corrected transfer. In particular original dilation is still the ordered finite exponential, not merely the semisimple eigenvalue.

\section{What the calculation establishes}
The new result is the exact remaining Gamma comparison on the unbounded subsequence, its original source/quotient maps, a closed evaluation of the free extra-node determinant, the finite signed mixed-Schur certificate, and a non-maximal-rank classification for the first mixed word required on that subsequence. The term in \eqref{eq:qlimit} has an actual nonzero $q$-scale limit for simple quartets. The mixed term in \eqref{eq:freeplusmixed}, however, remains necessary; neither its sign nor its growing-family cancellation has been supplied by the free determinant alone.

The unresolved expression after these calculations is explicit: the signed sum of $\log\det(I+X_N)$ and $\log\det(I-Y_N)$ in \eqref{eq:return}, or equivalently the mixed term in \eqref{eq:freeplusmixed}, combined with the inherited actual arithmetic $\delta_k^\sigma$. Formula \eqref{eq:cert} is a terminating fixed-input enclosure, not a tensor-uniform estimate. Formula \eqref{eq:arithmeticinterval} is its exact return to the original contradiction criterion. No claim of an RH contradiction, an actual counterexample, or failure of the absolute-base programme is made.

\paragraph{Evidence and source record.}
The accompanying scripts test the finite polynomial and covariance identities on explicitly declared exact calibrations. They do not assert that the calibration roots divide $2\xi$. The packet assumptions, DMR/QDC estimates, TWA/BRU/HC and existing STC/PR29 results retain their supplied scope. New execution logs, archive inventories and PDF-build receipts are separate from inherited records. Incoming archives and standalone files are retained byte-for-byte in the cumulative delivery. The new note has a complete independent proof body; retained historical source files are not silently edited.
\end{document}
